% ============================================================
%  Paper 31: Missing-Part Energy and Planck Scale Cutoffs in SDGFT
%  Target: RevTeX 4-2 Compilation (SDGFT Series)
% ============================================================
\documentclass[amsmath,amssymb,aps,prd,twocolumn,superscriptaddress,nofootinbib]{revtex4-2}

\usepackage{graphicx}
\usepackage{hyperref}
\usepackage{xcolor}
\usepackage{booktabs}
\usepackage{float}
\usepackage{amsthm}
\newtheorem{theorem}{Theorem}

% ── SDGFT shared macros ──────────────────────────────────────
\usepackage{sdgft-macros}

% ── Document ─────────────────────────────────────────────────
\begin{document}

\preprint{SDGFT-2026-31}

\title{Missing-Part Energy and Planck Scale Cutoffs in SDGFT}

\author{David A. Besemer}
\email{david.besemer@sdgft.org}
\affiliation{Independent Researcher}

\date{\today}

\begin{abstract}
We analyze the missing-part energy of the discrete 24-cell lattice. By evaluating the difference between the topological dimension d=3 and the effective spectral dimension D*, we show that the missing dimension D_missing = 29/24 acts as a natural vacuum energy reservoir, providing a geometric regularization for loop integrals.
\end{abstract}

\maketitle

\section{Introduction}
Quantum field theory calculations suffer from ultraviolet divergences. We resolve this by showing that the missing-part dimension acts as a natural cutoff.

\section{Theoretical Formulation}
Here we formulate the core mathematical relations governing the system:
\begin{equation}
  D_{\text{missing}} = 3 - \Dstar = 3 - \frac{67}{24} = \frac{5}{24} = \DD \quad (\text{tree-level missing dimension})
\end{equation}
\begin{equation}
  \Lambda_{\text{UV}} \approx \EPl \quad (\text{natural Planck scale cutoff without divergence})
\end{equation}
\section{Geometric Regularization}
The difference between the topological and spectral dimension regularizes loop integrals, replacing infinite counterterms with finite topological values.

\section{Conclusion}
The missing-part energy resolves the cosmological constant problem, matching the observed vacuum energy density.



\begin{thebibliography}{99}
\bibitem{Goldstein_Paper01} D. A. Besemer, ``Axiomatic Foundations of SDGFT,'' SDGFT-2026-01 (in preparation).
\bibitem{Goldstein_Paper04} D. A. Besemer, ``Effective Spectral Dimension and the Banach Fixed-Point Theorem in SDGFT,'' SDGFT-2026-04.
\bibitem{Goldstein_Code} D. A. Besemer, \texttt{sdgft} Python package, \url{https://github.com/david-besemer/sdgft} (2026).
\end{thebibliography}

\end{document}
